\documentclass[10pt,a4paper]{amsart}
\usepackage[margin=19mm]{geometry}
\usepackage{amsmath,amssymb,mathtools}
\usepackage[scr=rsfs,cal=euler]{mathalfa}
\usepackage{xfrac}
\usepackage[unicode,hidelinks,bookmarks=false]{hyperref}
\newcommand{\ie}{i.e.\ }
\newcommand{\cf}{cf.\ }
\newcommand{\resp}{respectively\ }
\newcommand{\Subsection}{\subsection}
\providecommand{\nobreakdash}{\nobreak\hbox{-}}
\setlength{\emergencystretch}{2em}
\multlinegap0pt
\usepackage{bm}
\usepackage{bbm}
\usepackage{dsfont}
\usepackage{xspace}
\usepackage{cancel}
\usepackage{mathtools}
\usepackage{amsthm}
\makeatletter
\@ifundefined{lemma}{\newtheorem{lemma}{Lemma}}{}
\makeatother
\newcommand{\bb}{\mathcal{B}}
\title[Strong chromatic index of bipartite graphs]{Strong chromatic index of bipartite graphs}
\author{Yanli Hao, Tianchi Yang,, Xingxing Yu}
\date{Source: arXiv:2606.23824v2 (CC BY 4.0)}
\begin{document}
\maketitle

\noindent\emph{Source and licence.} Strong chromatic index of bipartite graphs. Version arXiv:2606.23824v2 (CC BY 4.0), distributed under the Creative Commons Attribution 4.0 International licence, as recorded in the arXiv licence metadata for this exact version and shown on the abstract page https://arxiv.org/abs/2606.23824v2. Full rights notice: licenses/arxiv-2606.23824-CC-BY-4.0.txt.\par\smallskip

\noindent\textit{Editorial scope.} Counted unit: the Lemma whose source label is \texttt{clm BH} in the article (source0.tex lines 340--344, source label \texttt{clm BH}), delivered together with the article's own complete original proof of it (source0.tex lines 345--407, one \texttt{proof} environment of 63 lines). No mathematics is written, repaired or re-derived: statement and proof are the article's own text line for line. Required context retained uncounted: equ: Sx1 (lines 257--259). Every definition, display and label of those lines is retained unchanged. Omitted from this package and not redistributed: 1--256, 260--339, 408--524. Nothing inside a retained excerpt is omitted or altered: every retained body is delivered exactly as the article's lines are written, so no omission receipt of any kind accompanies this package. Citations: the retained excerpts carry no citation command at all, so no bibliography entry of the article is transcribed and no reference is redistributed. Reference adaptation: none was needed and none was made; every reference inside the delivered unit resolves inside it, so no unresolved reference or citation is produced. Printed slips kept literal: no printed word, symbol, hypothesis or number of the counted statement or of its proof was altered. Layout and class: the article's own class is \texttt{amsart} with packages including mystyle, tcolorbox, float; the wrapper is the lane's pinned \texttt{amsart} editorial wrapper at 10pt on A4, which restates the theorem-like environments the retained text needs and adds the article's own macro definitions that the retained text needs (\texttt{\textbackslash bb}), with extra wrapper packages (bm, bbm, dsfont, xspace, cancel, mathtools). The delivered numbering is the unit's own. Assets: no retained span contains a figure environment, a graphics-inclusion command, a PDF-page inclusion, a table or a TikZ picture; no image file is copied into this package, the package declares no asset dependency and no image file is redistributed with it. Licence: version arXiv:2606.23824v2 of this article is distributed under the Creative Commons Attribution 4.0 International licence, as recorded in the arXiv licence metadata for this exact version and shown on the abstract page https://arxiv.org/abs/2606.23824v2. A full rights notice is delivered at licenses/arxiv-2606.23824-CC-BY-4.0.txt. Characters: every retained excerpt is pure ASCII, so the delivered engine, XeLaTeX with its default Latin Modern text font, reports no missing character.
\begin{equation}\label{equ: Sx1}
    S_{A,1} =\sum_{x\in A_2}d_1(x)=|E_H(A_2,B_1)|=|B_1|\Delta_B-|E_H(A_1,B_1)|=\gamma \Delta_A\Delta_B
\end{equation}   

\begin{lemma}\label{clm BH} 
\[
  |B(H)| \ge    \frac{\Delta_B}{\Delta_A}S_{A,3}+\frac{\Delta_A}{\Delta_B} S_{B,3}- 3\Delta_B S_{A,2} -3\Delta_A S_{B,2} + 4\gamma \Delta_A^2\Delta_B^2
          +o(\Delta_A^2\Delta_B^2)   \]
\end{lemma}

\begin{proof} 
 Let $\bb_1(H)=\big\{abcd\in \bb(H): bc\notin E_H(A_2,B_2)\big\}$ and $\bb_2(H)=\big\{abcd\in \bb(H): bc\in E_H(A_2,B_2)\big\}$.  Then $|\bb(H)|=|\bb_1(H)|+|\bb_2(H)|$. We will bound $|\bb_1(H)|$ and $|\bb_2(H)|$. 

For $|\bb_1(H)|$, we count the ordered paths $abcd \in \mathcal{B}_1(H)$ based on the location of the vertex $c$. Recall that $ab\notin E_H(A_2,B_2)$ and $cd \in E_H(A_2, B_2)$. When $c \in A_2$, we have $d\in B_2$, $b\in B_1$, and $a\in A$. Hence, for each choice of $c\in A_2$, there are $d_1(c)d_2(c)$ choices for the ordered path $bcd$, and for each choice of such $bcd$, there are $\Delta_B -1$ choices for $a$ to form the ordered path $abcd$. Thus, the number of ordered paths $abcd$ in $\bb_1(H)$ with $c\in A_2$ is 
\begin{align*}
    \sum_{c\in A_2}d_1(c)d_2(c)(\Delta_B -1)
    &= \sum_{c\in A_2}d_1(c)(\Delta_A-d_1(c))(\Delta_B -1)\\
    &= (\Delta_B-1)\Delta_A  \sum_{c\in A_2}d_1(c)-(\Delta_B-1) \sum_{c\in A_2}d_1(c)^2\\
    &\ge  \Delta_A (\Delta_B-1) S_{A,1}- \Delta_B S_{A,2} \\
    &=%\overset{\eqref{equ: Sx1}}{=} 
    \gamma\Delta_A^2 \Delta_B^2  - \Delta_B S_{A,2} +o(\Delta_A^2 \Delta_B^2) \quad (\mbox{by } \eqref{equ: Sx1}).
\end{align*} 
Similarly, $c \in B_2$ implies $d\in A_2$, $b\in A_1$, and $a\in B$, and  the number of ordered paths $abcd$ in $\bb_1(H)$ with $c\in B_2$ is at least
 $\gamma\Delta_A^2 \Delta_B^2- \Delta_A S_{B,2}+o(\Delta_A^2 \Delta_B^2)$. Hence, 
%By relabeling $c$ to $x$ (when $c\in A_2$) or $y$ (when $c\in B_2$),  we have 
\[
    |\bb_1(H)|\ge 2\gamma\Delta_A^2 \Delta_B^2- \Delta_B S_{A,2}- \Delta_A S_{B,2}+o(\Delta_A^2 \Delta_B^2).
\] 
 
Next, we bound $|\bb_2(H)|$. For an ordered path $abcd\in \bb_2(H)$, we have  $bc,cd\in E_H(A_2,B_2)$ and $ab\notin E_H(A_2,B_2)$.
If $b\in A_2$ and $c\in B_2$ then $a\in B_1$ and $d\in A_2$; so for each such choice of $bc$, the number of  ordered paths $abcd$  in $\bb_2(H)$ with $b\in A_2$ and $c\in B_2$ is $d_1(b)(d_2(c)-1)=d_1(b)(\Delta_B-d_1(c)-1)$. If $b\in B_2$ and $c\in A_2$ then $a\in A_1$ and $d\in B_2$; so for each such choice of $bc$, the number of paths $abcd$ in $\bb_2(H)$ with $b\in B_2$ and  $c\in A_2$ is $d_1(b)(d_2(c)-1)=d_1(b)(\Delta_A-d_1(c)-1)$.
Hence, 
\[|\bb_2(H)|=\sum_{bc\in E(H), b\in A_2, c\in B_2}d_1(b)(\Delta_B-d_1(c)-1)+ \sum_{bc\in E(H), b\in B_2,c\in A_2}d_1(b)(\Delta_A-d_1(c)-1).\]
%Summing over all edges $xy \in E_H(A_2, B_2)$, we obtain
By renaming $b$ to $x$ and $c$ to $y$ (when $b\in A_2$) or $b$ to $y$ and $c$ to $x$ (when $b\in B_2$), we have  
    \begin{align*}\label{eqB2}
        |\bb_2(H)|
        &= \sum_{xy\in E(H), x\in A_2, y\in B_2}\Big(-2d_1(x)d_1(y)+ (\Delta_B -1 )d_1(x)+ (\Delta_A -1)d_1(y) \Big).
    \end{align*}
    Note that 
    $$\sum_{xy\in E(H), x\in A_2, y\in B_2} d_1(x)= \sum_{x\in A_2} d_1(x)d_2(x)=\sum_{x\in A_2}d_1(x)\Big(\Delta_A-d_1(x)\Big);$$
  %  \overset{\eqref{equ: Sx1}}{=} \gamma \Delta_A^2\Delta_B-S_{A,2}, $$ 
    hence
    \begin{align*}
     & \quad \sum_{xy\in E(H), x\in A_2,y\in B_2}(\Delta_B-1)d_1(x)\\
     &=(\Delta_B-1)\Big(\Delta_A\sum_{x\in A_2}d_1(x)-\sum_{x\in A_2}d_1(x)^2\Big)\\
     & =
     (\Delta_B -1 )( \gamma \Delta_A^2\Delta_B-S_{A,2}) \quad (\mbox{by } \eqref{equ: Sx1})\\
     & \ge \gamma \Delta_A^2\Delta_B^2-\Delta_B S_{A,2} -\gamma \Delta_A^2\Delta_B\\ &=\gamma \Delta_A^2\Delta_B^2-\Delta_B S_{A,2}+o(\Delta_A^2\Delta_B^2). 
   \end{align*}  
    Similarly, 
    $$\sum_{xy\in E(H),x\in A_2,y\in B_2}(\Delta_A-1) d_1(y)\ge  \gamma \Delta_A^2\Delta_B^2-\Delta_A S_{B,2}+o(\Delta_A^2\Delta_B^2) .$$
    Further note that 
     \begin{align*}
    &\quad \sum_{xy \in E(H),x\in A_2,y\in B_2} 2d_1(x)d_1(y) \\
    &=\Delta_A\Delta_B\sum_{xy \in E(H),x\in A_2,y\in B_2} 2\cdot \frac{d_1(x)}{\Delta_A}\cdot \frac{d_1(y)}{\Delta_B} \\   
    &\le \Delta_A\Delta_B\sum_{xy \in E(H),x\in A_2,y\in B_2}  \left(\left(\frac{d_1(x)}{\Delta_A}\right)^2+\left(\frac{d_1(y)}{\Delta_B}\right)^2\right)\\
    &=\frac{\Delta_B}{\Delta_A}\sum_{x\in A_2} d_1(x)^2(\Delta_A-d_1(x))+\frac{\Delta_A}{\Delta_B}\sum_{y\in B_2} d_1(y)^2(\Delta_B-d_1(y))\\
    &=-\frac{\Delta_B}{\Delta_A}\sum_{x\in A_2} d_1(x)^3 -\frac{\Delta_A}{\Delta_B} \sum_{y\in B_2} d_1(y)^3 + \Delta_B \sum_{x\in A_2} d_1(x)^2 + \Delta_A \sum_{y\in B_2} d_1(y)^2\\
     &=-\frac{\Delta_B}{\Delta_A}S_{A,3} -\frac{\Delta_A}{\Delta_B} S_{B,3} + \Delta_B S_{A,2} + \Delta_A S_{B,2} .
\end{align*}
Therefore,
    \begin{align*}
        |\bb_2(H)|&\ge    \frac{\Delta_B}{\Delta_A}S_{A,3}+\frac{\Delta_A}{\Delta_B} S_{B,3}- 2\Delta_B S_{A,2} -2\Delta_A S_{B,2} + 2\gamma \Delta_A^2\Delta_B^2
        +o(\Delta_A^2\Delta_B^2)
    \end{align*}
 
Hence, by combining the above bounds for $|\bb_1(H)|$ and $|\bb_2(H)|$, we have 
 \[
  |\bb_1(H)|+|\bb_2(H)|\ge  \frac{\Delta_B}{\Delta_A}S_{A,3}+\frac{\Delta_A}{\Delta_B} S_{B,3}- 3\Delta_B S_{A,2} -3\Delta_A S_{B,2} + 4\gamma \Delta_A^2\Delta_B^2
       +o(\Delta_A^2 \Delta_B^2). \]
Now the assertion of the lemma follows, since $|\bb(H)|=|\bb_1(H)|+|\bb_2(H)|$. 
\end{proof}

\end{document}
